% Generated draft: requires mathematical and source review.
\documentclass[11pt]{article}
\usepackage[margin=25mm]{geometry}
\usepackage{fontspec}
\setmainfont{Latin Modern Roman}
\usepackage{amsmath,amssymb,mathtools,unicode-math}
\setmathfont{Latin Modern Math}
\usepackage{enumitem,needspace,xurl,hyperref}
\setlength{\parindent}{0pt}
\setlength{\parskip}{5pt}
\setlength{\emergencystretch}{3em}
\newcommand{\N}{\mathbb{N}}
\newcommand{\Z}{\mathbb{Z}}
\newcommand{\Q}{\mathbb{Q}}
\newcommand{\R}{\mathbb{R}}
\newcommand{\C}{\mathbb{C}}
\newcommand{\F}{\mathbb{F}}
\newcommand{\A}{\mathbb{A}}
\newcommand{\PP}{\mathbb{P}}
\newcommand{\E}{\mathbb{E}}
\newcommand{\eps}{\varepsilon}
\newcommand{\dd}{\,\mathrm{d}}
\begin{document}
\section*{UnsolvedMath ID: 30001479}
\textbf{Source problem code:} catalogue entry
\subsection{Definitions and mathematical statement}
The supplied question concerns representations of a quiver with potential \((Q,S)\), an invariant denoted by \(E\), and mutation. The reference to negative simple representations indicates the usual decorated-representation interpretation, but the record does not explicitly specify this interpretation or the hypotheses required for mutation.

Under that interpretation, over a coefficient field \(k\), a decorated representation is a pair \(\mathcal M=(X,V)\), where \(X\) is a representation of the Jacobian algebra associated with \((Q,S)\) and \(V=(V_i)_{i\in Q_0}\) is a family of vector spaces indexed by the vertex set \(Q_0\). A representation is indecomposable if it is nonzero and cannot be expressed as a direct sum of two nonzero decorated representations. The negative simple representation \(\mathcal S_i^-\) at a vertex \(i\) has \(X=0\), \(V_i=k\), and \(V_j=0\) for \(j\ne i\). These definitions explain the conventional terminology; the field, finiteness conditions, and conventions for the Jacobian algebra remain unspecified in the supplied record. The symbol \(E(\mathcal M)\) denotes the source's \(E\)-invariant, whose definition and normalization are not supplied.

\textbf{Question (conditional formulation of the supplied statement).} For every quiver with potential \((Q,S)\) in the intended, unspecified class and every indecomposable \((Q,S)\)-representation \(\mathcal M\), does
\[
E(\mathcal M)=0
\quad\Longrightarrow\quad
\exists\,r\in\Z_{\geq 0},\ (Q^{(0)},S^{(0)}),\ i,\ k_1,\ldots,k_r
\quad
\mu_{k_r}\cdots\mu_{k_1}
\bigl(Q^{(0)},S^{(0)},\mathcal S_i^-\bigr)
\simeq (Q,S,\mathcal M)
\]
hold? Here \(i\) is a vertex of \(Q^{(0)}\), each \(k_j\) is a vertex at which the corresponding intermediate mutation is defined, and every mutation acts on both the quiver with potential and its decorated representation. For \(r=0\), the mutation sequence is empty. The symbol \(\simeq\) expresses identification of the terminal quiver with potential and representation with the prescribed data; its precise equivalence convention must be recovered from the source.

The displayed formulation makes explicit one conventional meaning of mutation generation. It does not resolve whether the source fixes the starting quiver with potential, what class of potentials is allowed, or which equivalence and mutation conventions are intended. In particular, nondegeneracy, reduction, restrictions on loops or oriented \(2\)-cycles, and hypotheses ensuring that successive mutations exist are not stated in the record. A fully specified unconditional statement requires those missing qualifications.
\subsection{Short English statement}
The supplied problem asks whether every indecomposable representation with vanishing \(E\)-invariant can be obtained from a negative simple representation by finitely many valid mutations. The decorated-representation interpretation is suggested by the terminology, but the ambient hypotheses, definition of \(E\), and precise mutation conventions require source clarification.
\subsection{Sources}
\begin{itemize}
\item Supplied local UnsolvedMath registry record, permanent ID \(30001479\), source problem code \texttt{catalogue entry}, titled \emph{Mutation Generation of E-Invariant-Zero Representations}, in the Oberwolfach Reports Open Problems collection. The record gives the proposed year as 2010 and names no individual proposer.
\item Supplied report citation: \emph{Interactions between Algebraic Geometry and Noncommutative Algebra} (2010), Oberwolfach Reports, pp.\ 1319--1381, DOI \texttt{10.4171/owr/2010/22}. \url{https://doi.org/10.4171/owr/2010/22}. The background describes Oberwolfach Reports as documenting mathematical workshops and their open-problem discussions; no problem-specific page is supplied.
\item Supplied literature source labeled Derksen--Weyman--Zelevinsky: \url{https://arxiv.org/abs/0904.0676}.
\item Supplied literature source labeled Plamondon: \url{https://arxiv.org/abs/1004.0830}.
\item The supplied literature assessment was checked on \texttt{2026-08-21}. It labels the problem open and states that the \(E\)-invariant has mutation and cluster-categorical characterizations while the exact mutation-generation assertion remains unresolved. This is supplied dated curation derived from the OWR record, not an independently verified assessment. The record explicitly requests independent verification of time-sensitive status and cited claims.
\item The registry labels the statement as corrected and verified and reports a manual semantic, typographic, and LaTeX review of its cleanup. These are supplied editorial qualifications; they do not establish the missing mathematical hypotheses or constitute independent mathematical review of this draft.
\end{itemize}
\end{document}
